Research on reproducing the canonical stylized facts of financial markets spans parametric return models, latent-state frameworks, multi-asset dependence modeling, and deep generative approaches.
This section reviews the key developments that motivate the present work, with particular attention to the Gaussian HMM limitation identified by \citet{ryden1998stylized} that our results challenge, and to the cross-asset dependence models that we use in our multi-asset extension.

\subsection{Stylized Facts of Asset Returns}

The Black--Scholes--Merton framework \cite{black1973pricing} assumes that log-returns are independently and identically distributed (i.i.d.) Gaussian random variables.
This assumption, while analytically convenient, is contradicted by a large body of empirical evidence accumulated over six decades.
\citet{mandelbrot1963variation} demonstrated that speculative price changes exhibit heavy tails (leptokurtosis), with extreme returns occurring far more frequently than a Gaussian model predicts.
\citet{fama1970efficient} confirmed negligible linear autocorrelation in raw returns, consistent with informational efficiency, but this does not imply independence: nonlinear dependence persists in the form of volatility clustering.
\citet{schwert1989does} documented that the autocorrelation of absolute returns decays slowly, often following a hyperbolic rather than exponential pattern, providing the empirical signature of clustered volatility.

\citet{cont2001empirical} systematized these observations into a canonical set of stylized facts that hold across asset classes, geographic markets, and sampling frequencies:
\begin{enumerate}
    \item \textit{Heavy tails}: the marginal distribution of returns exhibits positive excess kurtosis, with power-law tail decay.
    \item \textit{Absence of linear autocorrelation}: raw returns are approximately uncorrelated, consistent with market efficiency.
    \item \textit{Volatility clustering}: the autocorrelation of absolute (or squared) returns is positive and decays slowly, indicating that large price movements tend to be followed by large movements of either sign.
\end{enumerate}
Any generative model for financial time series must simultaneously reproduce all three facts to be considered statistically faithful.
Failure to capture even one, as we will show for the Gaussian i.i.d.\ generator (distributional failure) and GARCH(1,1) (distributional failure despite temporal fidelity), renders the synthetic data inadequate for downstream applications such as risk measurement and portfolio backtesting.

\subsection{Parametric Volatility Models}

The ARCH model of \citet{engle1982autoregressive} introduced conditional heteroscedasticity into time series econometrics, modeling the variance of returns as a function of past squared residuals.
The generalized ARCH (GARCH) extension of \citet{bollerslev1986generalized} added lagged conditional variance, producing the workhorse GARCH(1,1) specification
\begin{equation}
    \sigma_t^2 = \omega + \alpha\, r_{t-1}^2 + \beta\, \sigma_{t-1}^2
    \label{eq:garch_lit}
\end{equation}
where $\alpha + \beta < 1$ ensures stationarity.
The single-decay parameter $\beta$ directly controls the persistence of the conditional variance process, producing exponential ACF decay that closely approximates the empirical pattern when $\alpha + \beta$ is near unity.
However, GARCH models assume a unimodal innovation distribution and cannot represent discrete market regimes.
As we demonstrate empirically, this unimodal assumption leads to catastrophic distributional failure: GARCH(1,1) achieves only 13\% KS pass rate despite matching the observed ACF-MAE almost exactly (0.047 vs.\ 0.048 observed).

Parametric stochastic-volatility models treat log-variance as a latent diffusion, but parameter estimation requires computationally intensive methods such as particle filtering or Markov-chain Monte Carlo, and they are orthogonal to the synthetic-data-generation problem addressed here.

Jump-diffusion models \cite{merton1976option, kou2002jump} augment the geometric Brownian motion with Poisson-driven jumps in the price level, producing heavy-tailed marginal distributions.
\citet{merton1976option} used Gaussian jump sizes, while \citet{kou2002jump} introduced double-exponential jumps to capture the asymmetry between upside and downside tail risk.
These models generate heavy tails but lack a mechanism for volatility persistence: after a jump event, the volatility returns immediately to its baseline level, failing to reproduce the clustering phenomenon.

\subsection{Hidden Markov Models and the Ryd\'{e}n et al.\ Limitation}

\citet{hamilton1989new} introduced the Markov-switching model to economics, demonstrating that U.S.\ GNP growth is well described by a two-state regime-switching process.
The model posits a discrete latent state $S_t$ that evolves as a first-order Markov chain, with the observable variable drawn from a state-dependent distribution.
HMMs have since been applied to financial time series for regime identification \cite{kim1999has}, volatility estimation \cite{rossi2006volatility}, and trading strategy design \cite{nguyen2018hidden}.

The seminal evaluation of HMMs against stylized facts is due to \citet{ryden1998stylized}, who fitted Gaussian-emission HMMs with $K = 2$ and $K = 3$ states to daily Swedish equity returns.
They established three key findings.
First, HMMs produce heavy-tailed marginal distributions through the Gaussian mixture mechanism, even though each component is thin-tailed.
Second, HMMs produce negligible linear autocorrelation in returns, consistent with the observed stylized fact.
Third, HMMs \emph{fail} to reproduce the slow decay of the ACF of absolute returns (volatility clustering).
The third finding follows from a structural property of Markov chains: the sojourn time in any state follows a geometric distribution, which implies exponential decay of the ACF.
For $K = 2$--$3$ states, the exponential decay is too fast to approximate the hyperbolic-like persistence observed empirically.
This result has been widely cited as a fundamental limitation of standard HMMs for financial modeling.

\citet{bulla2006stylized} proposed hidden semi-Markov models (HSMMs) as a remedy.
HSMMs replace the geometric sojourn-time distribution with flexible parametric alternatives such as negative binomial or log-series, allowing arbitrarily slow ACF decay.
They demonstrated that HSMMs with $K = 2$--$4$ states and negative-binomial sojourn distributions successfully reproduce all three stylized facts, confirming that the Ryd\'{e}n limitation is fundamentally about the sojourn-time distribution, not the HMM framework itself.

The present work offers an alternative resolution: rather than modifying the sojourn-time distribution, we increase the state resolution.
With $K \geq 3$ states and quantile-based initialization, the Baum-Welch algorithm learns a transition matrix where high-volatility states exhibit sufficient self-transition probability (diagonal dominance) to sustain persistence.
The resulting mixture of $K$ geometric sojourn-time distributions, each with a different rate parameter, can approximate slow ACF decay, effectively achieving an HSMM-like effect within the standard HMM framework.
This insight has not been explored in the prior literature, likely because \citet{ryden1998stylized} only tested $K = 2$--$3$ and subsequent work adopted the HSMM path.

\subsection{Discrete HMM with Jump Augmentation}

\citet{alswaidan2026hybrid} addressed the Ryd\'{e}n limitation through a different mechanism: a Poisson jump-duration process that forces the model to dwell in tail states for durations drawn from a Poisson distribution, overriding the geometric sojourn time.
Their discrete framework used Laplace quantile binning to partition continuous observations into categorical states, estimated the transition matrix by direct frequency counting, and employed Student-$t$ within-bin emission distributions.
The jump mechanism introduced two hyperparameters ($\epsilon$ for jump probability, $\lambda$ for mean duration) calibrated via grid search.
The framework achieved 91--95\% KS pass rates across 424 US equities and partially reproduced volatility clustering.

The present work eliminates all three sources of complexity introduced by the discrete framework.
First, continuous Gaussian emissions replace Laplace quantile binning, removing discretization artifacts.
Second, the Baum-Welch algorithm jointly learns emissions and transitions, replacing frequency counting.
Third, no jump mechanism is needed at moderate $K$, removing two hyperparameters and the grid search entirely.

\subsection{Cross-Asset Dependence Models}
\label{sec:related_dep}

Generating jointly plausible multi-asset paths is a separate problem from univariate fidelity.
Two approaches dominate the quantitative finance literature.

\textit{Single Index Model (SIM).}
\citet{sharpe1963simplified} proposed a one-factor decomposition in which each asset's return is an affine function of a single market factor plus an idiosyncratic residual, $G_{i,t} = \alpha_i + \beta_i G_{M,t} + \eta_{i,t}$.
The SIM is parsimonious, analytically tractable, and is the workhorse reduction behind the mean-variance portfolio construction of \citet{markowitz1952portfolio}.
Its main limitations for synthetic data generation are that (i) the linear affine map distorts each asset's marginal distribution, so per-asset distributional fidelity degrades relative to the stand-alone marginal model, and (ii) it imposes a rank-one correlation structure in which all cross-asset dependence is mediated by the market factor.

\textit{Copulas.}
Sklar's theorem \cite{sklar1959fonctions, nelsen2006introduction} states that any joint cumulative distribution $H$ with continuous marginals $F_1, \dots, F_d$ factors uniquely as
\begin{equation}
    H(g_1, \dots, g_d) = C\bigl(F_1(g_1), \dots, F_d(g_d)\bigr),
\end{equation}
where the copula $C : [0,1]^d \to [0,1]$ is a multivariate distribution with uniform margins.
This factorization allows the joint distribution to be assembled from arbitrary per-asset marginals and any dependence structure $C$.
Within quantitative finance, the Gaussian and Student-t copulas \cite{demarta2005tcopula, mcneil2015quantitative} are standard choices because they admit analytic densities, are parameterized by a correlation matrix, and (for the Student-t copula) admit non-zero symmetric tail dependence controlled by the degrees-of-freedom parameter $\nu$.
Simulation via rank reordering \cite{iman1982distribution} exploits the uniqueness of Sklar's decomposition to inject the copula-sampled dependence structure onto independently simulated marginals without distorting any marginal distribution.

We adopt the copula construction of \cite{mcneil2015quantitative} for our cross-asset extension: per-asset CHMM simulations provide the marginals, Kendall's $\tau$ provides a robust rank-based correlation estimator, and the Student-t copula's degrees of freedom are selected by profile maximum likelihood.
The empirical ordering we report in Section~\ref{sec:cross_asset} (copulas ahead of SIM on both per-asset KS and cross-asset correlation reproduction) matches the ordering reported by \citet{alswaidan2026hybrid} for the discrete HMM marginals, suggesting that the copula advantage is intrinsic to the rank-reordering construction rather than specific to the marginal model.

\subsection{Deep Generative Models}

Deep generative models have attracted growing interest for financial time series synthesis.
\citet{takahashi2019modeling} applied generative adversarial networks (GANs) to stock return generation, finding that GAN-produced sequences visually resembled financial data but failed to reproduce ACF structure (the same volatility-clustering limitation as low-$K$ HMMs).
\citet{kwon2024can} systematically evaluated several GAN architectures against the Cont stylized facts and confirmed that volatility clustering remained the most difficult property to capture, unless the architecture was specifically designed to encode temporal dependence (for example, through recurrent discriminators).

\citet{alswaidan2026hybrid} benchmarked a GRU-based generator against HMM variants and found that it closed 83\% of the ACF-MAE gap relative to the HMM, but suffered variance collapse (0.6\% KS pass rate), indicating that distributional fidelity and temporal structure remain in tension for neural approaches.
The present study focuses on the parametric HMM approach, which offers interpretability (each state corresponds to a named market regime), computational efficiency ($O(TK^2)$ per Baum-Welch iteration), and guaranteed convergence (EM monotonicity), advantages that complement the flexibility of deep generative methods.

\subsection{Quality Assessment for Synthetic Time Series}

\citet{stenger2024thinking} surveyed the quality assessment landscape for synthetic time series, organizing evaluation metrics into distributional fidelity, temporal fidelity, and downstream utility categories.
Our seven-metric evaluation framework (KS and AD pass rates, excess kurtosis, ACF-MAE, Wasserstein-1, Hellinger distance, and quantile coverage) spans the first two categories, providing a comprehensive assessment that avoids reliance on any single metric.
As we will demonstrate, no single model dominates all seven metrics, and the choice of generator depends on which dimension of fidelity is prioritized.
\citet{assefa2020generating} made a complementary argument in the specific context of financial data, emphasizing that synthetic generators should be evaluated on their ability to preserve statistical properties relevant to the intended downstream application.
